%==============================================================================
\chapter{Périmètre du projet}
%==============================================================================

\section{Domaine couvert}
Le périmètre du présent cahier des charges couvre \textbf{la première partie}
du système : la gestion manuelle des opérations, leur planification et leur
suivi, ainsi que la préparation de l'architecture pour l'intégration future
des indicateurs automatisés.

\section{Modèle métier et règles de gestion}
\begin{xltabular}{\textwidth}{>{\bfseries}p{3cm} X}
\toprule
\textbf{Entité} & \textbf{Règles de gestion} \\
\midrule
\endhead
Fermier   & Peut posséder plusieurs fermes. \\
Ferme     & Peut comporter plusieurs sites d'apiculture. \\
Site      & Contient une ou plusieurs ruches, dispose d'une localisation géographique (latitude, longitude, altitude, date de mise en œuvre, éventuelles dates de déménagement / clôture). Un site peut accueillir des ruches de différents fermiers et de différentes fermes. \\
Ruche     & Composée d'un compartiment de base (socle / corps) et d'au plus \textbf{5} étages d'extension (hausses). \\
Corps / Hausse & Peut accueillir des cadres de cire installés sur des \emph{slots} identifiés. Le socle/corps doit contenir \textbf{au moins 1} cadre, la capacité maximale d'un corps comme d'une hausse est de \textbf{10 cadres}. \\
Cadre     & Occupe un emplacement (slot) unique dans le corps ou la hausse. \\
Agent apiculteur & Réalise les visites et inspections. Une ruche peut être suivie successivement par plusieurs agents, mais elle est sous la responsabilité d'\textbf{un seul agent} à un instant donné. \\
Agent superviseur & Approuve le planning de visite des ruches dont il a la charge. \\
\bottomrule
\end{xltabular}

\begin{encadre}[nobreak=true]
\textbf{Contrainte de composition :} 1 ruche = 1 socle/corps obligatoire + 0 à 5 hausses,
chaque corps/hausse = 1 à 10 cadres, un cadre par slot.
\end{encadre}

\begin{figure}[htbp]
\centering
\begin{tikzpicture}[
  every node/.style={font=\small},
  ent/.style={draw=mielfonce,fill=grisdoux,rounded corners=3pt,minimum height=9mm,inner xsep=7pt,thick},
  fl/.style={-{Latex[length=2mm]},mielfonce,thick}
]
\node[ent] (a) {Fermier};
\node[ent,right=7mm of a] (b) {Ferme};
\node[ent,right=7mm of b] (c) {Site};
\node[ent,right=7mm of c] (d) {Ruche};
\node[ent,below=9mm of d] (e) {Corps / Hausse};
\node[ent,left=7mm of e] (f) {Cadre};
\draw[fl] (a)--(b); \draw[fl] (b)--(c); \draw[fl] (c)--(d);
\draw[fl] (d)--(e); \draw[fl] (e)--(f);
\node[font=\scriptsize,mielfonce,above=0.5mm] at ($(a)!0.5!(b)$) {1..n};
\node[font=\scriptsize,mielfonce,above=0.5mm] at ($(b)!0.5!(c)$) {1..n};
\node[font=\scriptsize,mielfonce,above=0.5mm] at ($(c)!0.5!(d)$) {1..n};
\end{tikzpicture}
\caption{Hiérarchie des entités métier du système Zümm.}
\label{fig:hierarchie}
\end{figure}

\begin{figure}[htbp]
\centering
\begin{tikzpicture}[font=\small]
% périmètre du site
\draw[draw=mielfonce,thick,rounded corners=6pt,fill=grisdoux] (0,0) rectangle (11,4.4);
\node[anchor=north west,mielfonce,font=\bfseries] at (0.2,4.25) {Site apicole};
\node[anchor=north west,font=\footnotesize] at (0.2,3.7)
  {latitude / longitude / altitude, date de mise en oeuvre};
% trois ruches installées sur le site
\foreach \x in {0.9,4.3,7.7}{
  \begin{scope}[shift={(\x,0.7)}]
    \fill[miel!30,draw=mielfonce,thick] (0,0) rectangle (1.4,0.6);
    \fill[miel!40,draw=mielfonce,thick] (0,0.6) rectangle (1.4,1.2);
    \fill[mielfonce] (-0.12,1.2)--(1.52,1.2)--(0.7,1.7)--cycle;
    \draw[mielfonce,thick] (-0.1,-0.12) rectangle (1.5,0);
  \end{scope}
}
\node[font=\scriptsize] at (1.6,0.45) {Ruche};
\node[font=\scriptsize] at (5.0,0.45) {Ruche};
\node[font=\scriptsize] at (8.4,0.45) {Ruche};
% ressources externes
\node[anchor=west,font=\footnotesize,align=left] at (11.4,3.2)
  {Ressources\\(nectar / pollen)\\a environ 1 km};
\draw[-{Latex[length=2mm]},mielfonce,dashed] (11,3.0) -- (11.35,3.0);
\node[anchor=west,font=\footnotesize,align=left] at (11.4,1.2)
  {Distances de securite\\100 / 200 / 300 m\\3 a 4 km des cultures\\traitees};
\end{tikzpicture}
\caption{Schéma d'un site apicole : plusieurs ruches géolocalisées et contraintes d'emplacement.}
\label{fig:site}
\end{figure}

\section{Utilisateurs et droits d'accès}
\begin{xltabular}{\textwidth}{>{\bfseries}p{3.2cm} X}
\toprule
\textbf{Profil} & \textbf{Rôle et droits} \\
\midrule
\endhead
Fermier / Propriétaire & Consulte les tableaux de bord de production, de rentabilité et de synthèse, décide de la clôture d'un site ou du déplacement de ruches. \\
Agent apiculteur & Réalise les visites, remplit les rapports, exécute les opérations sur les ruches sous sa responsabilité. \\
Agent superviseur & Approuve les plannings de visite des ruches dont il est responsable. \\
Responsable / Manager & Surveille les alertes (chute de production ou d'effectif) et déclenche les inspections imminentes. \\
Administrateur & Gère les paramètres système, les unités de mesure, les seuils et les référentiels. \\
\bottomrule
\end{xltabular}

\section{Hors périmètre}
\begin{itemize}
  \item L'acquisition physique et le déploiement matériel des capteurs (projet ultérieur).
  \item L'automatisation temps réel de la collecte des indicateurs (interface prévue mais non implémentée dans cette phase).
\end{itemize}

\section{Perspectives d'évolution}
Au-delà du périmètre de cette phase, l'évolution suivante pourra être étudiée :
\begin{itemize}
  \item \textbf{Application mobile native (React Native), distincte de la PWA.}
  La solution retenue est une PWA, qui assure déjà la parité web / mobile et le
  fonctionnement hors-ligne (voir l'annexe des technologies). Une application
  native séparée ne serait engagée que si un besoin propre au natif l'imposait :
  notifications \emph{push} fiables sur iOS, géolocalisation en arrière-plan, ou
  distribution par les magasins d'applications. Elle réutiliserait le noyau métier
  TypeScript et l'API REST, mais constituerait une seconde base de code d'interface
  à maintenir --- à n'engager que sur exigence explicite du client.
\end{itemize}

